% concatenated from main.tex
%% ===================================================================
%% An inequality between the annihilation number and the Caro-Wei
%% sum, proving a TxGraffiti conjecture.
%%
%% Target: math.CO (arXiv), combinatorics note
%% ===================================================================
\documentclass[11pt]{article}
\usepackage[margin=1in]{geometry}
\usepackage{amsmath,amssymb,amsthm}
\usepackage[dvipsnames]{xcolor}
\usepackage[colorlinks=true,linkcolor=MidnightBlue,citecolor=MidnightBlue,
            urlcolor=MidnightBlue]{hyperref}
\usepackage{url}

% Theorem environments (each type has its own counter)
\newtheorem{theorem}{Theorem}
\newtheorem{corollary}{Corollary}
\newtheorem{conjecture}{Conjecture}
\theoremstyle{remark}
\newtheorem{remark}{Remark}

\newcommand{\degseq}{d_1 \le d_2 \le \cdots \le d_n}

\title{An annihilation-number Caro-Wei bound:\\
       a TxGraffiti conjecture and an independence-number bracket}
\author{Chakshu Gupta\\[2pt]
  {\small College of Computing, Georgia Institute of Technology}\\
  {\small \texttt{cgupta65@gatech.edu}}}
\date{}

\begin{document}
\maketitle

\begin{abstract}
Automated conjecturing programs scan collections of graphs for inequalities
between invariants that no stored graph violates, then offer the survivors for
proof or refutation. TxGraffiti, one such program, conjectured that every
nontrivial connected graph $G$ satisfies
$\alpha(G) \ge \bigl(a(G) + R(G)\bigr)/\Delta(G)$, where $\alpha$ is the
independence number, $a$ the annihilation number, $R$ the residue, and
$\Delta$ the maximum degree. Established only for two special families of
graphs, the conjecture has otherwise remained open. The note proves the
degree-sequence inequality $a \le \tfrac{\Delta+1}{2}W$, where $W$ is the
Caro-Wei sum; the same inequality is known for the independence number in
place of $a$. Combined with the classical lower bounds $\alpha \ge R$ and
$\alpha \ge W$, it proves the conjecture for every connected graph of maximum
degree at least three, and a direct argument settles maximum degree two; the
conjecture fails only for the single edge, of maximum degree one. The
inequality also brackets the independence number between the polynomial-time
quantities $R$ and $a$, within a factor $(\Delta+1)/2$. The
conjecture's bound is sharp, with equality attained, for instance, by the
complete graph on four vertices.
\end{abstract}

\section{Introduction}\label{sec:intro}

The Graffiti system~\cite{Fajtlowicz1988graffiti} makes conjectures by searching
a database of graphs for an inequality between invariants that no stored graph
violates, then offering the surviving relation as a conjecture. Its successor
TxGraffiti is the source of the following.

\begin{conjecture}[\cite{Davila2025reverie}]\label{conj:c1}
If $G$ is a nontrivial connected graph, then
$\alpha(G) \ge \bigl(a(G) + R(G)\bigr)/\Delta(G)$, and this bound is sharp.
\end{conjecture}

All graphs considered are finite, simple, and undirected. For a graph $G$ on $n$
vertices with $m$ edges and degree sequence $\degseq$ in nondecreasing order,
$\alpha(G)$ is the independence number; $\Delta(G) = d_n$ the maximum degree;
$a(G)$ the annihilation number, the largest integer $j$ such that
$d_1 + \cdots + d_j \le m$; and $R(G)$ the residue, the number of zeros
remaining when the Havel-Hakimi process~\cite{Havel1955, Hakimi1962} is iterated on the degree sequence until
it terminates. The independence number is NP-hard to compute~\cite{Karp1972},
whereas $a$ and $R$ are polynomial-time computable from the degree sequence.

The annihilation number is an upper bound on $\alpha$~\cite{Pepper2009}, while
the residue~\cite{Favaron1991} and the Caro-Wei sum $W$~\cite{Caro1979, Wei1981}
are lower bounds; $W$ sums $1/(d+1)$ over all vertices. The conjecture is open in
general, established only for regular bipartite and cubic K\"onig-Egerv\'ary
graphs~\cite{Davila2025reverie}. Existing bounds estimate $\alpha$
itself: the Degree Sequence Index Strategy bounds it from above
by annihilation-type quantities~\cite{CaroPepper2014}, and the Caro-Wei sum
estimates it within the best-possible factor $(\Delta+1)/2$, that is,
$\alpha \le \tfrac{\Delta+1}{2}W$~\cite{Boppana2018}. Proving the conjecture,
however, requires a bound on the annihilation number itself, which prior
estimates of $\alpha$ do not supply.

Theorem~\ref{thm:vehicle} provides such a bound: $a \le \tfrac{\Delta+1}{2}W$,
obtained by applying the argument of~\cite{Boppana2018} to the annihilation head
in place of a maximum independent set. Combined with the classical lower bounds,
it proves the conjecture for every connected graph of maximum degree at least
two (Theorem~\ref{thm:main}); the single edge, of maximum degree one, is the sole
exception, where the conjecture fails. The conjectured bound $(a+R)/\Delta$ is
sharp, attained for instance at the complete graph $K_4$
(Corollary~\ref{cor:sharp}), though for $\Delta \ge 3$, it is dominated by
$R$ (Corollary~\ref{cor:dominated}). The same inequality also brackets
the independence number between the polynomial-time quantities $R$ and
$a$, which differ by at most a factor $(\Delta+1)/2$ (Corollary~\ref{cor:approx}).

\section{Results}\label{sec:results}

For the nondecreasing degree sequence $\degseq$, the Caro-Wei sum is
$W = \sum_{i=1}^{n} \tfrac{1}{d_i+1}$, the index $i$ running over the
$n$ vertices. Three classical bounds on $\alpha$ are used throughout this
section: Favaron's residue bound $\alpha \ge R$~\cite[Thm.~3.3]{Favaron1991}, the
Caro-Wei bound $\alpha \ge W$~\cite{Caro1979, Wei1981}, and Pepper's
annihilation bound $\alpha \le a$~\cite{Pepper2009}. The residue dominates
the Caro-Wei sum, $R \ge W$, by Theorem~1.1 of~\cite{Favaron1991}.

\begin{theorem}\label{thm:vehicle}
Every graphic degree sequence with maximum degree $\Delta \ge 1$ satisfies
\[
  a \ \le\ \frac{\Delta+1}{2}\, W.
\]
\end{theorem}

\begin{proof}
Let $H = \{1, \ldots, a\}$ index the $a$ smallest degrees and
$T = \{a+1, \ldots, n\}$ the remaining $n-a$; call these the head and the tail.

\medskip\noindent\emph{Head degree-sum bound.}
The annihilation number $a$ is the largest $j$ with $d_1 + \cdots + d_j \le m$, so
the head satisfies $\sum_{i \in H} d_i \le m$. Each edge contributes $2$ to the
degree sum, hence $\sum_{i=1}^{n} d_i = 2m$; as the head and tail partition the
indices,
\[
  \sum_{i \in T} d_i \ =\ 2m - \sum_{i \in H} d_i \ \ge\ 2m - m \ =\ m
    \ \ge\ \sum_{i \in H} d_i .
\]
Every degree is at most $\Delta = d_n$, so the $n-a$ tail terms sum to at most
$\Delta(n-a)$. Chaining this with the previous line,
\begin{equation}\label{eq:headsum}
  \sum_{i \in H} d_i \ \le\ \sum_{i \in T} d_i \ \le\ \Delta\,(n-a).
\end{equation}

\medskip\noindent\emph{Pointwise inequality.}
For every integer $k$ with $0 \le k \le \Delta$,
\begin{equation}\label{eq:pointwise}
  \frac{1}{k+1} + \frac{k}{\Delta(\Delta+1)} \ \ge\ \frac{2}{\Delta+1}.
\end{equation}
Multiplying through by the positive number $\Delta(\Delta+1)(k+1)$ and collecting
terms,~\eqref{eq:pointwise} is equivalent to
\[
  \Delta(\Delta+1) + k(k+1) - 2\Delta(k+1) \ \ge\ 0,
\]
whose left-hand side equals $(\Delta-k)(\Delta-k-1)$. Since $0 \le k \le \Delta$,
the integer $\Delta - k$ is nonnegative: if $\Delta - k = 0$, the product is
$0 \cdot (-1) = 0$; if $\Delta - k \ge 1$, then $\Delta - k - 1 \ge 0$ and the
product is the product of two nonnegative integers. In either case it is
nonnegative, with equality precisely when one factor vanishes, that is, at
$k = \Delta$ and $k = \Delta - 1$.

\medskip\noindent\emph{Summation.}
Split $W$ over the head and tail. On the tail every degree is at most $\Delta$, so
$\tfrac{1}{d_i+1} \ge \tfrac{1}{\Delta+1}$, and the $n-a$ tail terms contribute at
least $(n-a)/(\Delta+1)$. On the head apply~\eqref{eq:pointwise} with $k = d_i$,
legitimate as each $d_i \le \Delta$:
\[
  W \ \ge\ \sum_{i \in H}\!\left(\frac{2}{\Delta+1} - \frac{d_i}{\Delta(\Delta+1)}\right)
        + \frac{n-a}{\Delta+1}
    \ =\ \frac{2a}{\Delta+1}
        - \frac{\sum_{i \in H} d_i}{\Delta(\Delta+1)}
        + \frac{n-a}{\Delta+1}.
\]
By~\eqref{eq:headsum}, $\sum_{i \in H} d_i \le \Delta(n-a)$, so the middle term is
at least $-(n-a)/(\Delta+1)$; it cancels the tail contribution, leaving
$W \ge 2a/(\Delta+1)$. Rearranging gives $a \le \tfrac{\Delta+1}{2}W$.
\end{proof}

The hypothesis $\Delta \ge 1$ is necessary. With $\Delta = 0$, the edgeless
sequence on $n$ vertices has $a = n$ and $W = n$, so $\tfrac{\Delta+1}{2}W = n/2
< a$. The constant $(\Delta+1)/2$ is best possible. For every $\Delta \ge 1$,
any $\Delta$-regular sequence of even order $n$ has $a = n/2$ and
$W = n/(\Delta+1)$, so $a = \tfrac{\Delta+1}{2}W$.

\begin{theorem}[Resolution of Conjecture~\ref{conj:c1}]\label{thm:main}
Every connected graph $G$ with $\Delta(G) \ge 2$ satisfies
\[
  \Delta(G)\,\alpha(G) \ \ge\ a(G) + R(G).
\]
\end{theorem}

\begin{proof}
The hypothesis $\Delta \ge 2$ excludes $K_2$, the only connected graph with
maximum degree $1$, for which $\Delta\,\alpha = 1 < a + R = 2$.
If $\Delta = 2$, then $G$ is a path or a cycle, and in both cases $a = \alpha$
by a direct calculation from the degree sequence. With $\alpha \ge R$,
\[
  a + R \ \le\ 2\alpha \ =\ \Delta\,\alpha.
\]
If $\Delta \ge 3$, then $\tfrac{\Delta+1}{2} \le \Delta-1$. By
Theorem~\ref{thm:vehicle} and the Caro-Wei bound $\alpha \ge W$,
\[
  a \ \le\ \tfrac{\Delta+1}{2}\,W \ \le\ \tfrac{\Delta+1}{2}\,\alpha
    \ \le\ (\Delta-1)\,\alpha.
\]
Adding $\alpha \ge R$ gives $\Delta\,\alpha = (\Delta-1)\alpha
+ \alpha \ge a + R$.
\end{proof}

\begin{corollary}[Sharpness]\label{cor:sharp}
For every connected graph $G$ with $\Delta(G) \ge 2$, the bound of
Theorem~\ref{thm:main} holds with equality iff $(\Delta-1)\alpha = a$
and $\alpha = R$. The complete graph $K_4$ attains equality, with
$\Delta\,\alpha = a + R = 3$.
\end{corollary}

\begin{proof}
By the proof of Theorem~\ref{thm:main}, $(\Delta-1)\alpha \ge a$ for every
$\Delta \ge 2$; Favaron's bound gives $\alpha \ge R$. Hence
$\Delta\,\alpha = (\Delta-1)\alpha + \alpha \ge a + R$, with equality iff
both summand inequalities are equalities, that is, $(\Delta-1)\alpha = a$ and
$\alpha = R$. The complete graph $K_4$ has $\alpha = 1$, $a = 2$, $R = 1$, and
$\Delta = 3$, so both equalities hold and $\Delta\,\alpha = a + R = 3$.
\end{proof}

\begin{corollary}[Domination by $R$]\label{cor:dominated}
For every connected graph $G$ with $\Delta(G) \ge 3$,
\[
  \frac{a(G) + R(G)}{\Delta(G)} \ \le\ R(G).
\]
\end{corollary}

\begin{proof}
For $\Delta \ge 3$, $\tfrac{\Delta+3}{2} \le \Delta$. Theorem~\ref{thm:vehicle}
and $W \le R$ give $a \le \tfrac{\Delta+1}{2}R$. Hence
\[
  a + R \ \le\ \tfrac{\Delta+1}{2}R + R
        \ =\ \tfrac{\Delta+3}{2}R \ \le\ \Delta R,
\]
and dividing by $\Delta$ gives the inequality.
\end{proof}

\begin{corollary}[Bracketing the independence number]\label{cor:approx}
For every connected graph with maximum degree $\Delta \ge 1$,
\[
  R \ \le\ \alpha \ \le\ a \ \le\ \tfrac{\Delta+1}{2}R.
\]
The outer ratio is sharp, attained by $K_{\Delta+1}$ for every odd $\Delta \ge 1$.
\end{corollary}

\begin{proof}
The bounds $\alpha \ge R$ and $\alpha \le a$ give $R \le \alpha \le a$;
Theorem~\ref{thm:vehicle} together with $W \le R$ extends this to
$a \le \tfrac{\Delta+1}{2}W \le \tfrac{\Delta+1}{2}R$. For every odd $\Delta \ge 1$, the complete
graph $K_{\Delta+1}$ has $\alpha = 1$, $a = (\Delta+1)/2$, and $R = W = 1$, so
$a = \tfrac{\Delta+1}{2}R$.
\end{proof}

\begin{remark}[Refining the Caro-Wei bracket]
The bracket $R \le \alpha \le a$ refines the Caro-Wei bracket
$W \le \alpha \le \tfrac{\Delta+1}{2}W$~\cite{Boppana2018}, since
$R \ge W$ tightens the lower endpoint and
Theorem~\ref{thm:vehicle} caps the upper endpoint at $\tfrac{\Delta+1}{2}W$.
Both endpoints are computable in polynomial time. The worst-case ratio of
upper to lower endpoint, $(\Delta+1)/2$, is unchanged, attained at
$K_{\Delta+1}$ for odd $\Delta$ (Corollary~\ref{cor:approx}).
\end{remark}

Both theorems were verified by exhaustive computation in exact
arithmetic.\footnote{\url{https://github.com/ChakshuGupta13/lab}}
Theorem~\ref{thm:vehicle} holds for all $3{,}166{,}851$ graphic degree
sequences of order at most $14$ with minimum degree at least $1$, and
Theorem~\ref{thm:main} holds for all $11{,}989{,}762$ connected graphs of
order $3$ through $10$, matching OEIS~A001349.\footnote{\url{https://oeis.org/A001349}}
Graphs were generated with \texttt{geng} from \texttt{nauty}~\cite{McKay2014nauty}.

\section{Conclusion}\label{sec:disc}

Previously the conjecture $\alpha \ge (a+R)/\Delta$ had been verified only for
regular bipartite graphs and cubic K\"onig-Egerv\'ary graphs. This note
resolves it for every connected graph of maximum degree at least two; the
single edge, with $\Delta = 1$, is the unique counter-example. The argument
reduces to the inequality $a \le \tfrac{\Delta+1}{2}W$, capping the ratio
between the annihilation number and the Caro-Wei sum at $(\Delta+1)/2$; the
constant is best possible, attained on regular sequences of even order.
Combined with the classical lower bounds $\alpha \ge W$ and $\alpha \ge R$,
this inequality forces the conjecture.

The conjectured bound $(a+R)/\Delta$ is itself dominated by $R$ for
$\Delta \ge 3$, so the conjecture names a relation among classical
degree-sequence quantities of $\alpha$ rather than tightening any of them.
The bracket $R \le \alpha \le a$ localizes $\alpha$ in polynomial
time within a factor $(\Delta+1)/2$, refining the Caro-Wei bracket. Equality
in the conjecture forces $\alpha = R$ and $(\Delta-1)\alpha = a$; the
complete graph $K_4$ attains both, as do paths and cycles with $R = \alpha$.
A structural description of all equality graphs is left open.

\bibliographystyle{alpha}
\bibliography{references}

\end{document}
